\section{Generación de la familia Bendersky--Glaisher--Kinkelin mediante regularización zeta}

Sea \(\mathcal Z_3(s):=\zeta(s)\) la realizaci\'on escalar del funcional
zeta tri\'adico en la normalizaci\'on de este volumen,
\(H_k=\sum_{j=1}^k j^{-1}\) (con \(H_0=0\)) y \(B_{k+1}\) el n\'umero de
Bernoulli de convenio \(B_1=-1/2\). Definimos

\begin{equation}
A_k=\exp\!\left(
\frac{H_kB_{k+1}}{k+1}-\mathcal Z_3'(-k)
\right),\qquad k\ge0.
\label{eq:bendersky}
\end{equation}

Los primeros niveles incluyen

\[
A_0=\sqrt{2\pi},
\]

y la constante de Glaisher--Kinkelin en \(A_1\). Para \(n\ge1\),

\begin{equation}
\zeta(2n+1)=(-1)^{n+1}
\frac{2(2\pi)^{2n}}{(2n)!}\log A_{2n}.
\label{eq:odd-zeta-bendersky}
\end{equation}

La constante de Ap\'ery, \(\zeta(3)\), es el primer nivel par no trivial.
Su presencia en el gas fot\'onico del \cref{ch:metrology} constituye una
realizaci\'on f\'isica posterior de este valor espectral.

La primera parte de la torre puede verse de forma sinóptica:
\begin{center}
\small
\begin{tabular}{@{}cll@{}}
\toprule
\(k\) & identidad estructural & valor terminal\\
\midrule
0 & \(A_0=\exp[-\zeta'(0)]=\sqrt{2\pi}\)
  & \(2.5066282746310005024\ldots\)\\
1 & \(\log A_1=1/12-\zeta'(-1)\)
  & \(1.2824271291006226369\ldots\)\\
2 & \(\log A_2=-\zeta'(-2)=\zeta(3)/(4\pi^2)\)
  & \(1.0309167521973921\ldots\)\\
3 & \(\log A_3=-11/720-\zeta'(-3)\)
  & \(0.9795555269428446\ldots\)\\
4 & \(\log A_4=-\zeta'(-4)\)
  & \(0.9920479745250403\ldots\)\\
\bottomrule
\end{tabular}
\end{center}

\begin{proof}[Relaci\'on con los valores impares]
La ecuaci\'on funcional de \(\zeta\), desarrollada en \(s=-2n\), da
\begin{equation}
\zeta'(-2n)=(-1)^n
\frac{(2n)!}{2(2\pi)^{2n}}\zeta(2n+1).
\label{eq:zeta-derivative-even-negative}
\end{equation}
Como \(B_{2n+1}=0\) para \(n\ge1\), la definici\'on
\cref{eq:bendersky} reduce \(\log A_{2n}=-\zeta'(-2n)\). Despejar
\(\zeta(2n+1)\) prueba \cref{eq:odd-zeta-bendersky}. El \`indice \(2n\)
conserva el nivel de derivaci\'on: no es una etiqueta editorial.
\end{proof}

